% Master CV source. Every version of the CV is built from this file.
%
% The version is set at compile time through \cvtarget (see the Makefile):
%   detailed        Detailed CV (the default)
%   academic        Academic CV
%   academic_twopager        Two-page academic CV
%   epidemiologist           Senior Epidemiologist roles
%   epidemiologist_twopager  Two-page CV for Senior Epidemiologist roles
%   datascientist            Data Scientist roles
%   datascientist_twopager   Two-page CV for Data Scientist roles
%
% \cvversion{<versions>}{<content>} shows <content> only in the listed versions.
% \ifcvversion{<versions>}{<content>}{<other content>} picks between two variants.

\documentclass[11pt,a4paper,sans]{moderncv} % Font sizes: 10, 11, or 12; paper sizes: a4paper, letterpaper, a5paper, legalpaper, executivepaper or landscape; font families: sans or roman
\moderncvstyle{classic} % CV theme - options include: 'casual' (default), 'classic', 'oldstyle' and 'banking'
\moderncvcolor{blue} % CV color - options include: 'blue' (default), 'orange', 'green', 'red', 'purple', 'grey' and 'black'

\usepackage{ragged2e}
\usepackage[scale=0.85]{geometry} % Reduce document margins
\setlength{\hintscolumnwidth}{4cm} % Uncomment to change the width of the dates column
\usepackage{enumitem}
\usepackage[T1]{fontenc}
\usepackage{lmodern}

\usepackage{makecell}
\usepackage{multicol}
% Start a new page if less than #1 of space is left (keeps headings with their first entry)
% Subsection headings: bold, with a thin rule to the end of the line and more space above
\AtBeginDocument{%
  \renewcommand*{\subsectionfont}{\large\bfseries\upshape}%
  \renewcommand*{\subsectionstyle}[1]{{\subsectionfont\textcolor{subsectioncolor}{#1}}\hspace{0.6em}{\color{subsectioncolor!35}\leaders\hrule height 0.6pt depth 0pt\hfill\kern0pt}}%
}
\makeatletter
\RenewDocumentCommand{\subsection}{sm}{%
  \par\addvspace{2ex}%
  \phantomsection{}%
  \addcontentsline{toc}{subsection}{#2}%
  \cvitem[0ex]{}{\strut\subsectionstyle{#2}}%
  \par\nobreak\addvspace{1ex}\@afterheading}
\makeatother
\makeatletter
\newcommand{\needspace}[1]{\par\begingroup\setlength{\dimen@}{#1}\dimen@ii\pagegoal\advance\dimen@ii-\pagetotal
  \ifdim\dimen@>\dimen@ii\ifdim\dimen@ii>\z@\vfil\fi\penalty-\@M\fi\endgroup}
\makeatother
\renewcommand\theadalign{bc}
\renewcommand\theadfont{\bfseries}
\renewcommand\theadgape{\Gape[4pt]}
\renewcommand\cellgape{\Gape[4pt]}

% CV versions
\providecommand{\cvtarget}{detailed}
\ExplSyntaxOn
% Versions are compared as strings, so underscores in names match however they were read
\str_new:N \g__cv_target_str
\str_gset:Ne \g__cv_target_str { \cvtarget }
\bool_new:N \l__cv_match_bool
\cs_new_protected:Npn \__cv_match:n #1
  {
    \bool_set_false:N \l__cv_match_bool
    \clist_map_inline:nn {#1}
      { \str_if_eq:eeT { \tl_to_str:n {##1} } { \g__cv_target_str } { \bool_set_true:N \l__cv_match_bool \clist_map_break: } }
  }
\__cv_match:n { detailed, academic, academic_twopager, epidemiologist, epidemiologist_twopager, datascientist, datascientist_twopager }
\bool_if:NF \l__cv_match_bool
  { \PackageError { CV } { Unknown~CV~version~`\g__cv_target_str' } { Set~\token_to_str:N \cvtarget \ to~one~of~the~versions~listed~at~the~top~of~the~file. } }
\NewDocumentCommand { \cvversion } { m +m }
  { \__cv_match:n {#1} \bool_if:NT \l__cv_match_bool {#2} }
\NewDocumentCommand { \ifcvversion } { m +m +m }
  { \__cv_match:n {#1} \bool_if:NTF \l__cv_match_bool {#2} {#3} }
\ExplSyntaxOff

% Two-pagers only keep a heading with the first few lines below it, so sections fill the page
\let\cvneedspace\needspace
\ifcvversion{academic_twopager,epidemiologist_twopager,datascientist_twopager}{%
  \renewcommand{\needspace}[1]{\cvneedspace{4\baselineskip}}}{}

\firstname{James} % Your first name
\familyname{Mba Azam} % Your last name

\title{Curriculum Vitae}
\address{London School of Hygiene \& Tropical Medicine}{Keppel Street, London WC1E 7HT, United Kingdom}
\email{james.azam@lshtm.ac.uk}
\homepage{https://jamesmbaazam.github.io/}{}
\social[orcid]{0000-0001-5782-7330}

% Sections that appear in a different place in some versions
\newcommand{\cvtechnicalcompetencies}{%
\cvitem{Programming tools}{R, Stan, Bash, GNU Make, GitHub \& Git}
\cvitem{Machine learning}{Random forests and other tree-based methods, and gradient boosting (XGBoost, CatBoost, LightGBM), using frameworks including \texttt{tidymodels}}
\cvitem{Forecast evaluation}{\texttt{scoringutils}}
\cvitem{Dashboard tools}{RShiny, Quarto Dashboards}
\cvitem{Markup languages}{LaTeX, Markdown, R Markdown, HTML}
% \cvitem{Operating Systems}{MacOS, Linux/Ubuntu, Microsoft Windows}
% \cvitem{Presentation}{Quarto Presentations, MS PowerPoint, LaTeX Beamer}
\cvitem{Presentation}{Quarto Presentations, LaTeX Beamer}
\cvitem{Project management}{Agile/scrum}
% \cvitem{Graphic Design}{Inkscape, Adobe suite, Camtasia, Canva, CapCut}
}

\newcommand{\cvresearchsoftware}{%
\needspace{8\baselineskip}
\section{Research Software}
\cvversion{detailed,academic,datascientist}{%
\cventry{Nov, 2022 -- Present}{Developer \& maintainer}{\texttt{epichains} \& \texttt{bpmodels}}{}{}{R packages for simulating and analysing transmission chains using branching process models. \texttt{epichains} supersedes \texttt{bpmodels} and has over 6,000 CRAN downloads.\newline Citation: \textbf{Azam, J.M.}, Funk, S. \& Finger, F. (2024). \texttt{epichains}: Simulating and Analysing Transmission Chain Statistics Using Branching Process Models. R package version 0.1.1. \url{https://doi.org/10.5281/zenodo.13945039}}
\cventry{Nov, 2022 -- Aug, 2025}{Co-maintainer}{\texttt{EpiNow2}}{}{}{R package for estimating real-time infectious disease cases and time-varying epidemiological parameters, with over 62,000 CRAN downloads.\newline Citation: Abbott, S., Hellewell, J., Sherratt, K., Gostic, K., Hickson, J., Badr, H.S., DeWitt, M., \textbf{Azam, J.M.}, Lison, A. \& Funk, S. (2026). \texttt{EpiNow2}: Estimate and Forecast Real-Time Infection Dynamics. R package version 1.9.0. \url{https://doi.org/10.5281/zenodo.3957489}}
}%
\cvversion{epidemiologist,datascientist_twopager}{%
\cventry{Nov, 2022 -- Present}{Developer \& maintainer}{\texttt{epichains} \& \texttt{bpmodels}}{}{}{R packages for simulating and analysing transmission chains using branching process models. \texttt{epichains} supersedes \texttt{bpmodels} and has over 6,000 CRAN downloads.}
\cventry{Nov, 2022 -- Aug, 2025}{Co-maintainer}{\texttt{EpiNow2}}{}{}{R package for estimating real-time infectious disease cases and time-varying epidemiological parameters, with over 62,000 CRAN downloads.}
}%
\cvversion{detailed,academic,epidemiologist,datascientist}{\cventry{2023 -- Present}{Contributor}{\texttt{epinowcast}, \texttt{primarycensored}, and \texttt{scoringutils}}{}{}{R packages for nowcasting (\texttt{epinowcast}), estimating primary event censored delay distributions (\texttt{primarycensored}), and evaluating forecasts (\texttt{scoringutils}).}}
}

\begin{document}

\makecvtitle % Print the CV title

\cvversion{detailed,academic,epidemiologist,datascientist}{%
\section{Professional Summary}
\cvversion{detailed,academic}{%
\cvitem{}{Infectious disease modeller developing mathematical, statistical and machine learning methods to evaluate outbreak response interventions and predict how outbreaks spread. As a Research Fellow at the Centre for Mathematical Modelling of Infectious Diseases, London School of Hygiene \& Tropical Medicine, currently developing spatial statistical and machine learning models to predict the timing and location of vaccine-derived poliovirus outbreaks in Africa and the Middle East, as part of a Gates Foundation-funded project. First-author research on outbreak response vaccination published in \textit{Epidemics} and \textit{Vaccine}. Developer of the \texttt{epichains} R package for analysing transmission chains, and former co-maintainer of the widely used \texttt{EpiNow2} R package for real-time outbreak analytics. Committed to teaching and capacity building, with experience teaching on the MSc Health Data Science, leading the \textit{Modelling for Pandemic Preparedness and Response} short course for participants from more than 15 African countries, and supervising and mentoring interns and early career researchers. Completed Advance HE's Diversifying Leadership Programme and chairs LSHTM's Black Early Career Staff Network.}
}%
\cvversion{epidemiologist}{%
\cvitem{}{Infectious disease epidemiologist and modeller with a PhD in Applied Mathematics and experience delivering analyses for outbreak response and preparedness across the UK, Africa and Latin America. Has provided modelling decision support to South Africa's National Advisory Group on Immunization on the threat of polio re-emergence, contributed to WHO consultations and seminars on pandemic H5N1 influenza and mpox analytics, and currently analyses multi-pathogen wastewater surveillance data with the National Institute for Communicable Diseases and partners. Experienced in leading capacity building, including leading the \textit{Modelling for Pandemic Preparedness and Response} short course for over 65 participants from more than 15 African countries, and in leading a staff network as its elected chair, supported by Advance HE's Diversifying Leadership Programme.}
}%
\cvversion{datascientist}{%
\cvitem{}{Health data scientist with a PhD in Applied Mathematics and a track record of building machine learning models, statistical models and reproducible analysis pipelines for public health. Currently developing tree-based and gradient boosting models to predict the spatial spread of vaccine-derived poliovirus outbreaks at the London School of Hygiene \& Tropical Medicine, and previously co-maintained the \texttt{EpiNow2} R package (over 62,000 CRAN downloads) as a Research Software Engineer. Skilled in R, Stan, \texttt{tidymodels}, Git and GNU Make, and in turning complex analyses into open-source tools, dashboards and training for non-specialist users.}
}%
}%

\cvversion{datascientist,datascientist_twopager}{%
\needspace{6\baselineskip}
\section{Technical Competencies}
\cvtechnicalcompetencies
}%

\needspace{6\baselineskip}
\section{Education}
\cventry{Feb, 2018 -- Apr, 2022}{PhD Applied Mathematics}{The South African Centre for Epidemiological Modelling and Analysis (SACEMA),
Stellenbosch University}{South Africa}{}{}
\cvitem{Thesis title}{\emph{Modelling outbreak response intervention strategies for decision-making}}
\cvversion{detailed,academic}{%
\cvitem{Supervisors}{Prof. Juliet R.C. Pulliam \& Prof. Matthew J. Ferrari}
\vspace{0.4cm}
}%

\cventry{Mar, 2016 -- Mar, 2018}{MSc Mathematics}{The South African Centre for Epidemiological Modelling and Analysis (SACEMA), Stellenbosch University}{South Africa}{}{}
\cvversion{detailed,academic,epidemiologist,datascientist}{\cvitem{Thesis title}{\emph{Assessment of a measles outbreak response vaccination campaign, and two measles parameter estimation methods}}}
\cvversion{detailed,academic}{%
\cvitem{Supervisor}{Prof. Juliet R.C. Pulliam}
\vspace{0.4cm}
}%

\cventry{Aug, 2014 -- Jun, 2015}{Structured MSc Mathematical Sciences}{African Institute for Mathematical Sciences (AIMS)}{Senegal}{}{}
\cvversion{detailed,academic,epidemiologist,datascientist}{\cvitem{Thesis title}{\emph{Hedging Conditional Value at Risk}}}
\cvversion{detailed,academic}{%
\cvitem{Supervisor}{Dr. Maciej J. Capinski}
\vspace{0.4cm}
}%

\cventry{Aug, 2009 -- Jun, 2013}{BSc Actuarial Science}{Kwame Nkrumah University of Science and Technology}{Ghana}{}{}
\cvversion{detailed,academic,epidemiologist,datascientist}{\cvitem{Thesis title}{\emph{Predicting optimal periods for the sale of cocoa in Ghana for the year 2013, using futures contracts}}}
\cvversion{detailed,academic}{%
\cvitem{Supervisor}{Mr. Vincent K. Dedu}
\vspace{0.4cm}
}%

\ifcvversion{academic_twopager,epidemiologist_twopager,datascientist_twopager}{\needspace{6\baselineskip}}{\needspace{25\baselineskip}}
\section{Professional Experience}
\cvversion{detailed,academic,epidemiologist,datascientist}{%
\cventry{Sep, 2025 -- Present}{Research Fellow}{\textit{Centre for Mathematical Modelling of Infectious Diseases}, London School of Hygiene \& Tropical Medicine}{London}{United Kingdom}{
	\begin{itemize}
		\ifcvversion{datascientist}{%
		\item Work on Objective 1 of the Gates Foundation-funded project \textit{Modeling to Support Polio Eradication and OPV Cessation}: develop spatial statistical and machine learning models, including tree-based and gradient boosting models, that incorporate novel data sources to predict the timing and location of type 2 vaccine-derived poliovirus outbreaks and outbreak responses in Africa and the Middle East. First-author paper in preparation for \textit{PLOS Global Public Health}.
		}{%
		\item Work on Objective 1 of the Gates Foundation-funded project \textit{Modeling to Support Polio Eradication and OPV Cessation}: develop spatial statistical and machine learning models that incorporate novel data sources to predict the timing and location of type 2 vaccine-derived poliovirus outbreaks and outbreak responses in Africa and the Middle East. First-author paper in preparation for \textit{PLOS Global Public Health}.
		}%
		% \item Supervise master's student projects. (Uncomment once supervising.)
		\item Teach on the Data Challenge module of the MSc Health Data Science: facilitate the Git and GitHub session and the R for data exploration and visualisation sessions, and tutor a group of 4--5 students on data challenges set by external organisations, including the UK Health Security Agency (UKHSA), UNICEF, private start-ups, and public health data science companies.
	\end{itemize}
}
\cventry{Nov, 2022 -- Aug, 2025}{Research Software Engineer}{\textit{Epiverse-TRACE}, London School of Hygiene \& Tropical Medicine}{London}{United Kingdom}{
	\begin{itemize}
		\item Developed and maintained the \texttt{epichains} R package (superseding \texttt{bpmodels}) for simulating and analysing transmission chains, with over 6,000 CRAN downloads.
		\item Co-maintained the \texttt{EpiNow2} R package for estimating the time-varying reproduction number ($R_t$) and forecasting infections in real time. \texttt{EpiNow2} is widely adopted by researchers and public health organisations, including the US CDC, WHO, and ministries of health such as the Philippines', in responses to outbreaks of COVID-19 and measles, and to the mpox outbreaks in the US in 2022 and in the Democratic Republic of the Congo. It has served as a basis and benchmark for newer methods and software for real-time $R_t$ estimation and forecasting, including \texttt{epinowcast}, EpiFusion, ConvRt, \texttt{estimateR}, and the FASSSTER platform.
		\item Conducted infectious disease modelling research for outbreak response and preparedness, including an analysis of forecast performance using daily and weekly aggregated data (first-author paper in preparation for \textit{PLOS Global Public Health}).
		\cvversion{detailed,academic,datascientist}{%
		\item Developed technical and non-technical resources such as R package vignettes, function references, tutorials, how-to guides, and blog posts, for example for \texttt{EpiNow2}, \texttt{bpmodels}, and \texttt{epichains}.
		}%
		\item Provided and supported outbreak analytics and data science trainings through Epiverse-TRACE, for example at the KEMRI-Wellcome Trust Research Programme in Kilifi, Kenya, and the Epiverse-TRACE pre-summit training at the MRC Unit The Gambia at LSHTM.
		\cvversion{detailed,academic,datascientist}{%
		\item Supported outbreak analytics open-source communities, including epiforecasts and epinowcast, through code contributions, code reviews, seminars, and training activities.
		}%
	\end{itemize}
}
}%
\cvversion{academic_twopager}{%
\cventry{Sep, 2025 -- Present}{Research Fellow}{\textit{Centre for Mathematical Modelling of Infectious Diseases}, London School of Hygiene \& Tropical Medicine}{London}{United Kingdom}{Spatial statistical and machine learning models to predict vaccine-derived poliovirus outbreaks in Africa and the Middle East (Gates Foundation-funded). Teaching on the MSc Health Data Science.}
\cventry{Nov, 2022 -- Aug, 2025}{Research Software Engineer}{\textit{Epiverse-TRACE}, London School of Hygiene \& Tropical Medicine}{London}{United Kingdom}{Developed the \texttt{epichains} R package (over 6,000 CRAN downloads) and co-maintained \texttt{EpiNow2} (over 62,000 CRAN downloads) for real-time outbreak analytics; delivered outbreak analytics training in Kenya and The Gambia.}
}%
\cvversion{epidemiologist_twopager}{%
\cventry{Sep, 2025 -- Present}{Research Fellow}{\textit{Centre for Mathematical Modelling of Infectious Diseases}, London School of Hygiene \& Tropical Medicine}{London}{United Kingdom}{Spatial statistical and machine learning models to predict the timing and location of vaccine-derived poliovirus outbreaks and outbreak responses in Africa and the Middle East (Gates Foundation-funded).}
\cventry{Nov, 2022 -- Aug, 2025}{Research Software Engineer}{\textit{Epiverse-TRACE}, London School of Hygiene \& Tropical Medicine}{London}{United Kingdom}{Co-maintained \texttt{EpiNow2}, used by the US CDC, WHO and ministries of health in responses to outbreaks of COVID-19, measles and mpox; developed the \texttt{epichains} R package for analysing transmission chains; delivered outbreak analytics training in Kenya and The Gambia.}
}%
\cvversion{datascientist_twopager}{%
\cventry{Sep, 2025 -- Present}{Research Fellow}{\textit{Centre for Mathematical Modelling of Infectious Diseases}, London School of Hygiene \& Tropical Medicine}{London}{United Kingdom}{Spatial statistical and machine learning models, including tree-based and gradient boosting models, to predict vaccine-derived poliovirus outbreaks in Africa and the Middle East (Gates Foundation-funded). Teach Git, GitHub and R on the MSc Health Data Science.}
\cventry{Nov, 2022 -- Aug, 2025}{Research Software Engineer}{\textit{Epiverse-TRACE}, London School of Hygiene \& Tropical Medicine}{London}{United Kingdom}{Developed the \texttt{epichains} R package and co-maintained \texttt{EpiNow2}; wrote package documentation, tutorials and how-to guides, and contributed code and code reviews to open-source outbreak analytics software.}
}%

\cvversion{detailed,academic,epidemiologist}{%
\cventry{Mar, 2025 -- Jul, 2025}{Technical Reviewer and Video Content Creator}{Introduction to Infectious Disease Modelling for Public Health}{Epiverse-TRACE online course}{}{
	\begin{itemize}
		\item Reviewed and edited English translations of outbreak analytics training materials originally developed in Spanish, ensuring clarity, technical accuracy, and alignment with the African context, including using examples from Africa.
		\item Created English versions of video explainers on modelling topics such as SIR models, Zika compartmental modelling, and model simulation in R, originally created in Spanish.
	\end{itemize}
}
}%

\ifcvversion{academic_twopager,epidemiologist_twopager,datascientist_twopager}{%
\cvversion{academic_twopager}{\cventry{Feb, 2022 -- Nov, 2022}{Junior Researcher}{The South African Centre for Epidemiological Modelling and Analysis (SACEMA)}{Stellenbosch}{South Africa}{Co-developed a polio transmission model to support South Africa's National Advisory Group on Immunization, and an R and GNU Make pipeline for detecting changes in reported COVID-19 severity.}}%
\cvversion{epidemiologist_twopager}{\cventry{Feb, 2022 -- Nov, 2022}{Junior Researcher}{The South African Centre for Epidemiological Modelling and Analysis (SACEMA)}{Stellenbosch}{South Africa}{Co-developed a spatial, age-structured polio transmission model to provide decision support to South Africa's National Advisory Group on Immunization on possible polio re-emergence.}}%
\cvversion{datascientist_twopager}{\cventry{Feb, 2022 -- Nov, 2022}{Junior Researcher}{The South African Centre for Epidemiological Modelling and Analysis (SACEMA)}{Stellenbosch}{South Africa}{Co-developed an R and GNU Make pipeline for detecting changes in the trends of reported COVID-19 severity in South Africa, and a spatial, age-structured polio transmission model.}}%
}{%
\cventry{Feb, 2022 -- Nov, 2022}{Junior Researcher}{The South African Centre for Epidemiological Modelling and Analysis (SACEMA)}{Stellenbosch}{South Africa}{
	\begin{itemize}
	 \item Co-developed a spatial, age-structured compartmental model of polio transmission in South Africa to provide analytics decision support to the country's National Advisory Group on Immunization in preparation for possible polio re-emergence, following wild poliovirus detections in 2022 in southern African countries such as Mozambique and Malawi.
		\item Co-developed an analysis pipeline (using \textit{R} and \textit{GNU make}) for detecting changes in the trends of reported COVID-19 severity in South Africa.
		\item Developed a model to analyse strategies for mitigating a hypothetical pathogen in the face of emerging variants.
		\item Contributed to the activities of SACEMA Modelling and Analytics Response Team (SMART).
	\end{itemize}
}
}%

\cvversion{detailed,academic,epidemiologist,datascientist}{%
\cventry{Feb, 2021 -- Sep, 2021}{Data Science Consultant (Part-time)}{Winning Insights and Models for Marvellous Yields (Wimmy)}{Cape Town}{South Africa}{
	\begin{itemize}
		\item Extended an existing COVID-19 model to analyse the impact of vaccine rollout in South Africa, Chile, Peru, and Zimbabwe.
		\item Contributed to a process to improve the calibration, validation, and run-time of the COVID-19 model.
		\cvversion{detailed,academic,epidemiologist}{%
		\item Led and contributed to LinkedIn articles about various health issues of immediate concern, including the potential impact of emerging COVID-19 variants, and addressing the issue of blood clotting following COVID-19 vaccination.
		}%
		\cvversion{detailed,academic,datascientist}{%
		\item Drafted guidelines for code testing and review within Wimmy.
		}%
		\item Contributed to the design of proposals for new analyses for Anglo American, including a study on Noise-Induced Hearing Loss among miners.
\end{itemize}}
}%
\cvversion{epidemiologist_twopager}{\cventry{Feb, 2021 -- Sep, 2021}{Data Science Consultant (Part-time)}{Winning Insights and Models for Marvellous Yields (Wimmy)}{Cape Town}{South Africa}{Extended a COVID-19 model to analyse the impact of vaccine rollout in South Africa, Chile, Peru, and Zimbabwe.}}%
\cvversion{datascientist_twopager}{\cventry{Feb, 2021 -- Sep, 2021}{Data Science Consultant (Part-time)}{Winning Insights and Models for Marvellous Yields (Wimmy)}{Cape Town}{South Africa}{Improved the calibration, validation, and run-time of a COVID-19 model.}}%

\cvversion{datascientist,datascientist_twopager}{\cvresearchsoftware}

\cvversion{detailed,academic,epidemiologist}{%
\needspace{8\baselineskip}
\section{Research Affiliations and Collaborations}
\cventry{Apr, 2026 -- Present}{Measles Analytics Hub}{Imperial College London}{United Kingdom}{}{Mentor on the Hub's mentoring scheme, supervising a project to estimate subnational measles immunity in Uganda.}
\cventry{Dec, 2023 -- Present}{German-West African Centre for Global Health and Pandemic Prevention (G-WAC)}{Kwame Nkrumah University of Science and Technology, Ghana, and Technische Universität Berlin}{Germany}{}{Co-designed the \textit{Modelling for Pandemic Preparedness and Response} short course at a curriculum development workshop (Dec, 2023) and have led it since, with three iterations (2024--2026) for over 65 participants from more than 15 African countries. Honorary Senior Research Fellow since Jan, 2025, in recognition of contributions to G-WAC's regional capacity building in infectious disease modelling and global health research. First-author paper on the course's development and delivery submitted to \textit{Epidemics}.}
\cventry{Jan, 2023 -- Present}{The South African Centre for Epidemiological Modelling and Analysis (SACEMA)}{Stellenbosch University}{South Africa}{}{Ongoing collaboration on research, training, and data sharing through honorary appointments, in recognition of sustained contributions to SACEMA's research and training activities, as Research Associate (Jan, 2023 -- Jul, 2026) and Research Fellow (Jul, 2026 -- Present).}
}%

\needspace{10\baselineskip}
\ifcvversion{detailed,academic}{\section{Publications}}{\section{Selected Publications}}

\cvversion{detailed,academic,epidemiologist,datascientist}{%
\needspace{6\baselineskip}
\subsection{Published: first author}
}%

\cvlistitem{\textbf{Azam, J.M.}, Pang, X., Are, E.B., Pulliam, J.R.C. \& Ferrari, M.J. (2023). Modelling outbreak response impact in human vaccine-preventable diseases: A systematic review of differences in practices between collaboration types before COVID-19. \textit{Epidemics}, 45, 100720. \url{https://doi.org/10.1016/j.epidem.2023.100720}}

\cvlistitem{\textbf{Azam, J.M.}, Saitta, B., Bonner, K., Ferrari, M.J. \& Pulliam, J.R.C. (2021). Modelling the relative benefits of using the measles vaccine outside cold chain for outbreak response. \textit{Vaccine}, 39(40), 5845--5853. \url{https://doi.org/10.1016/j.vaccine.2021.08.053}}

\cvlistitem{\textbf{Azam, J.M.}, Are, E.B., Pang, X., Ferrari, M.J. \& Pulliam, J.R.C. (2020). Outbreak response intervention models of vaccine-preventable diseases in humans and foot-and-mouth disease in livestock: a protocol for a systematic review. \textit{BMJ Open}, 10(10), e036172. \url{https://doi.org/10.1136/bmjopen-2019-036172}}

\cvversion{detailed,academic,epidemiologist,datascientist}{%
\needspace{6\baselineskip}
\subsection{Published: co-author}
}%

\cvversion{detailed,academic,epidemiologist}{%
\cvlistitem{Nightingale, E.S., Robert, A., Auzenbergs, M., \textbf{Azam, J.M.}, Villabona-Arenas, J., Fu, H., Sherratt, K., Funk, S., Jit, M., Atkins, K.E. \& O'Reilly, K. (2026). Bridging data and decision-making: A mathematical modelling perspective on infectious disease elimination. \textit{PLOS Global Public Health} (in press).}

}%
\cvversion{detailed,academic,epidemiologist,datascientist}{%
\cvlistitem{Cuartero, C.T., Lambert, J.W., Cuomo-Dannenburg, G., Fitzner, J., Ndumbi, P., Rice, C., Saulo, D., Waddell, L., \textbf{Azam, J.M.}, Doohan, P., Hicks, J.T., McCain, K., Morgenstern, C., Naidoo, T.M., Ruybal-Pesántez, S., Rawson, T., Santoni, C., Unwin, H.J.T., Kucharski, A., Cori, A., Bhatia, S. \& McCabe, R. (2026). The impact of ambiguously reported epidemiological parameters for infectious disease modelling and recommended best practices. \textit{Epidemics}, 100942. \url{https://doi.org/10.1016/j.epidem.2026.100942}}

}%

\cvversion{detailed,academic,academic_twopager,epidemiologist,epidemiologist_twopager}{%
\cvlistitem{van Zandvoort, K., Procter, S.R., \textbf{Azam, J.M.}, Sherratt, K. \& Davies, N.G. (2026). The risk of global Ebola virus spread is low: epidemiology of Ebola disease cases outside Africa, 1976 to May 2026. \textit{Eurosurveillance}, 31(24), 2600508. \url{https://doi.org/10.2807/1560-7917.ES.2026.31.24.2600508}}

}%
\cvversion{detailed,academic,academic_twopager,epidemiologist}{%
\cvlistitem{Aheto, J.M.K., Auzenbergs, M., Ferrari, M.J., Portnoy, A., Utazi, C.E., Kakaï, R.G., Gayawan, E., \textbf{Azam, J.M.} \& Nonvignon, J. (2026). Rebalancing power in infectious disease modelling: Toward inclusive and contextual approaches. \textit{PLOS Global Public Health}, 6(4), e0006220. \url{https://doi.org/10.1371/journal.pgph.0006220}}

}%
\cvversion{detailed,academic,epidemiologist,epidemiologist_twopager}{%
\cvlistitem{Ward, J., Lambert, J.W., Russell, T.W., \textbf{Azam, J.M.}, Kucharski, A.J., Funk, S., \dots \& Edmunds, W.J. (2025). Estimates of epidemiological parameters for H5N1 influenza in humans: a rapid review. \textit{BMC Infectious Diseases}, 25(1), 1755. \url{https://doi.org/10.1186/s12879-025-11933-z}}

}%
\cvversion{detailed,academic}{%
\cvlistitem{Kent, B.A., Holman, C., Amoako, E., Antonietti, A., \textbf{Azam, J.M.}, Ballhausen, H., \dots \& Weissgerber, T.L. (2022). Recommendations for empowering early career researchers to improve research culture and practice. \textit{PLOS Biology}, 20(7), e3001680. \url{https://doi.org/10.1371/journal.pbio.3001680}}

\cvlistitem{Czachur, M.V., Todd, M., Loureiro, T.G., \textbf{Azam, J.M.}, Nyeleka, S., Alblas, A. \& Davies, S.J. (2021). Opening the floor for discussion: A perspective on how scholars perceive attitudes to science in policymaking in South Africa. \textit{South African Journal of Science}, 117(1/2), 7948. \url{https://doi.org/10.17159/sajs.2021/7948}}

\cvlistitem{Lukwa, A.T., Siya, A., Zablon, K.N., \textbf{Azam, J.M.} \& Alaba, O.A. (2020). Socioeconomic inequalities in food insecurity and malnutrition among under-five children: within and between-group inequalities in Zimbabwe. \textit{BMC Public Health}, 20(1), 1199. \url{https://doi.org/10.1186/s12889-020-09295-z}}

}%
\cvversion{detailed,academic,epidemiologist}{%
\needspace{6\baselineskip}
\subsection{Published: Global Burden of Disease consortium}

\cvlistitem{Schumacher, A.E., Zheng, P., Barber, R.M., \dots, \textbf{Azam, J.M.}, \dots \& Alam, M. (2025). Global age-sex-specific all-cause mortality and life expectancy estimates for 204 countries and territories and 660 subnational locations, 1950--2023: a demographic analysis for the Global Burden of Disease Study 2023. \textit{The Lancet}, 406(10513), 1731--1810. \url{https://doi.org/10.1016/S0140-6736(25)01330-3}}

\cvlistitem{Naghavi, M., Kyu, H.H., Aalipour, M.A., \dots, \textbf{Azam, J.M.}, \dots \& Al Homsi, A. (2025). Global burden of 292 causes of death in 204 countries and territories and 660 subnational locations, 1990--2023: a systematic analysis for the Global Burden of Disease Study 2023. \textit{The Lancet}, 406(10513), 1811--1872. \url{https://doi.org/10.1016/S0140-6736(25)01917-8}}

\cvlistitem{GBD 2021 Antimicrobial Resistance Collaborators (incl. \textbf{Azam, J.M.}) (2024). Global burden of bacterial antimicrobial resistance 1990--2021: a systematic analysis with forecasts to 2050. \textit{The Lancet}, 404(10459), 1199--1226. \url{https://doi.org/10.1016/S0140-6736(24)01867-1}}

}%
\cvversion{detailed,academic,epidemiologist,datascientist}{%
\needspace{6\baselineskip}
\subsection{In review}

}%
\cvversion{detailed,academic,epidemiologist}{%
\cvlistitem{\textbf{Azam et al.}, Modelling for Pandemic Preparedness and Response Modular Shortcourse - Development and Delivery in Ghana. \textit{Submitted to Epidemics}}

}%
\cvversion{detailed,academic,epidemiologist,datascientist}{%
\cvlistitem{Korir et al., Replicability and Reusability of Epiverse-TRACE’s Tutorials on Outbreak Analytics and Applied Modelling in R. \textit{Submitted to Epidemics}}

}%
\cvversion{detailed,academic,epidemiologist}{%
\cvlistitem{Valle-Campos et al., Collaborative Curriculum Development and Delivery of Epidemiological Data Science Training in Africa. \textit{Submitted to Epidemics}}

\cvlistitem{Gómez-Bermeo et al., The Epi-Training Kit Translation and Localisation for Africa: Strengthening Capacity in Epidemiological Modelling through South–South Collaboration. \textit{Submitted to Epidemics}}

}%
\cvversion{detailed,academic,epidemiologist,datascientist}{%
\needspace{6\baselineskip}
\subsection{In preparation}

\cvlistitem{\textbf{Azam et al.}, Predicting the spatial spread of vaccine-derived polio outbreaks in Africa using machine learning. \textit{Target journal: PLOS Global Public Health}}

\cvlistitem{\textbf{Azam et al.}, Daily versus weekly data in epidemiological forecasting: computational and practical implications. \textit{Target journal: PLOS Global Public Health}}

}%
\cvversion{detailed,academic,datascientist}{%
\cvlistitem{Brand et al., A scalable marginalisation approach for double interval censored epidemiological delays.}

}%
\cvversion{detailed}{%
\needspace{6\baselineskip}
\subsection{Media}
\cvlistitem{\raggedright Forbes.com. Will Coronavirus Measures Slow The Spread Of Measles? Ask This Ghanaian. Jul 2020. URL: \url{https://www.forbes.com/sites/andrewwight/2020/07/29/will-coronavirus-measures-slow-measles-ask-this-ghanaian/}}

}%
\cvversion{detailed,academic,epidemiologist}{\cvresearchsoftware}

\cvversion{detailed,academic}{%
\needspace{6\baselineskip}
\section{Fellowships, Honours, and Awards}
\cventry{Sep, 2020 \& Sep, 2021}{Best student oral presentation}{SACEMA Research Days}{}{}{}
\cventry{Dec, 2019}{Stellenbosch University Library Open Access Fund}{\pounds 405}{}{}{A competitively awarded subsidy towards an Open Access Publication fee.}
}%
\cvversion{detailed,academic}{%
\cventry{Oct, 2018}{Big Data Policy Fellowship}{\pounds 4,000}{Development in Africa with Radio Astronomy (DARA)}{University of Manchester, United Kingdom}{A one-month fellowship that provided training in science and policy communication in the context of Big Data projects and in particular, the Square Kilometre Array in South Africa and Australia.}

}%
\cvversion{detailed,academic}{%
\cventry{Feb, 2018}{PhD Bursary}{R140,000}{The South African Centre for Epidemiological Modelling and Analysis (SACEMA)}{Stellenbosch University, South Africa}{Awarded competitively through the PhD application process.}

\cventry{Aug, 2017}{Best student oral presentation}{SACEMA Research Days}{}{}{}

\cventry{Feb, 2016}{Master's Bursary}{R120,000}{SACEMA}{Stellenbosch University, South Africa}{Awarded competitively through the master's application process.}

}%
% Grant Applications section hidden: the only entry was not awarded. Uncomment to show.
% \needspace{8\baselineskip}
% \section{Grant Applications}
% \cventry{Jun, 2026}{Contributor to proposal writing}{Gates Foundation Grand Challenges}{\textit{Estimating the Global Burden of Diarrheal Diseases} call}{}{\textit{Empirical Estimation of Pathogen-Specific Diarrheal Mortality and Morbidity Burden in Ghanaian Children Under Five: Integrating Health System Surveillance, Clinical Data, and Molecular Diagnostics}. Proposed budget: US\$1,253,500.}
% Outcome: not awarded.

\ifcvversion{detailed,academic,epidemiologist}{\needspace{14\baselineskip}}{\needspace{6\baselineskip}}
\section{Leadership}
\ifcvversion{detailed,academic,epidemiologist}{%
\cventry{Feb, 2023 -- Present}{Chair}{Black Early Career Staff Network}{London School of Hygiene \& Tropical Medicine}{London, United Kingdom}{
	\begin{itemize}
		\item \textbf{Elected Chair} (Dec, 2025 -- Present).
		\item \textbf{Acting Chair} (Apr, 2025 -- Dec, 2025): led the network through its formal recognition as an LSHTM staff network in Jun, 2025.
		\item \textbf{Lead, Formation Task Force} (Jul, 2024 -- Jun, 2025): led the case for formalising the network. Drafted the terms of reference, the cover sheet for LSHTM's EDI Committee, and the equality impact assessment; consulted members and incorporated their feedback; and advocated for formalisation at senior meetings, including with the EDI Office and LSHTM's EDI Committee. This work formed part of LSHTM's application for the Advance HE Race Equality Charter Bronze Award, which recognises institutional progress towards race equity.
		\item \textbf{Co-organiser} (Feb, 2023 -- Jul, 2024): co-organised the network informally before its formal recognition.
	\end{itemize}
}
}{%
\cventry{Feb, 2023 -- Present}{Chair}{Black Early Career Staff Network}{London School of Hygiene \& Tropical Medicine}{London, United Kingdom}{Elected Chair since Dec, 2025. Led the network through its formal recognition as an LSHTM staff network in Jun, 2025, work that formed part of LSHTM's application for the Advance HE Race Equality Charter Bronze Award.}
}%
\cvversion{detailed,academic}{%
\cventry{Jul, 2020 -- Jul, 2021}{President}{Association of Ghanaian Students in Stellenbosch University}{Stellenbosch University}{South Africa}{}
\cventry{Jun, 2019 -- Jun, 2020}{Vice President}{Association of Ghanaian Students in Stellenbosch University}{Stellenbosch University}{South Africa}{}
}%
%\cventry{2017}{Student Representative}{SACEMA Research Days planning committee}{Stellenbosch University}{South Africa}{}
%\cventry{2017}{Student Representative}{SACEMA Strategic Planning Meeting}{Stellenbosch University}{South Africa}{}
%\cventry{2009 -- 2013}{Member of several committees}{Pax Romana Local and Pax Romana Choir}{Kwame Nkrumah University of Science and Technology}{Ghana}{}

\cvversion{detailed,academic}{%
\needspace{10\baselineskip}
\section{Academic Service}
\needspace{4\baselineskip}
\subsection{Peer review}

\cvlistdoubleitem{\textit{Journal of Open Source Software} (2 R package papers)}{\textit{Bulletin of Mathematical Biology} (1)}
\cvlistdoubleitem{\textit{PLOS Global Public Health} (1)}{\textit{Scientific Reports} (1)}

\needspace{6\baselineskip}
\subsection{Equity, diversity and inclusion}
\cventry{Apr, 2025 -- Present}{Member}{Equity, Diversity and Inclusion Committee}{London School of Hygiene \& Tropical Medicine}{London, United Kingdom}{}
\cventry{Dec, 2022 -- Oct, 2023}{Active Member}{Equity, Diversity, Inclusion, and Accessibility Working Group}{Society of Research Software Engineering}{United Kingdom}{}
\cventry{Jul, 2023}{Career Panel Member}{Health Data Research UK Black Internship Programme}{London School of Hygiene \& Tropical Medicine}{United Kingdom}{}
}%

\needspace{10\baselineskip}
\ifcvversion{detailed,academic,epidemiologist,datascientist}{\section{Presentations}}{\section{Selected Presentations}}
\cvversion{detailed,academic,epidemiologist,datascientist}{%
\needspace{6\baselineskip}
\subsection{Conference Oral Presentations}

}%
\cvversion{detailed,academic,epidemiologist,datascientist,academic_twopager,datascientist_twopager}{%
\cvitem{Apr, 2025}{\texttt{epichains}: an R package for analysing transmission chains data. \textit{Humanitarian Modelling Dialogue Workshop, WHO Pandemic Hub, Berlin, Germany}}

\cvitem{Nov, 2024}{Benchmarking time-varying reproduction number, $R_t$, and forecasting models -- a case study of the \texttt{EpiNow2} R package. \textit{Infectious Disease Modelling (IDM) Conference, Bangkok, Thailand}}

}%
\cvversion{detailed,academic,datascientist,datascientist_twopager}{%
\cvitem{Jul, 2023}{Superseding R packages – lessons learned from transitioning to the \texttt{epichains} R package from \texttt{bpmodels}. \textit{Research Software London Southeast conference, Imperial College London, United Kingdom}}

}%
\cvversion{detailed,academic,epidemiologist}{%
\cvitem{Nov, 2021}{\textbf{JM Azam}, MJ Ferrari, JRC Pulliam. Modelling the relative benefits of using the measles vaccine outside cold chain for outbreak response. \textit{Epidemics 8, Elsevier (Virtual presentation)}}

}%
\cvversion{detailed,academic}{%
\cvitem{Sep, 2021}{\textbf{JM Azam}, X Pang, EB Are, MJ Ferrari, JRC Pulliam. Outbreak response intervention models of vaccine-preventable diseases in humans (1970--2019) -- A systematic review. \textit{SACEMA Research Days (PhD Student Presentation)}. Best student oral presentation.}

\cvitem{Sep, 2020}{\textbf{JM Azam}, MJ Ferrari, JRC Pulliam. Modelling the relative benefits of using the measles vaccine outside cold chain for outbreak response. \textit{SACEMA Research Days (PhD Student Presentation)}. Best student oral presentation.}

\cvitem{Oct, 2019}{\textbf{JM Azam}, MJ Ferrari, JRC Pulliam. Measles outbreak response: is it worthwhile to use the vaccine outside of the cold chain? \textit{Modelling in the Context of African Health (MoCAH)}, Durban, South Africa.}

\cvitem{Sep, 2019}{\textbf{JM Azam}, X Pang, EB Are, MJ Ferrari, JRC Pulliam. Outbreak response models of human vaccine preventable diseases (1970-2019) - A systematic review. \textit{SACEMA Research Days (PhD Student Presentation)}, Stellenbosch, South Africa.}

\cvitem{Aug, 2017}{\textbf{JM Azam}, JRC Pulliam. Measles in a school in Stellenbosch, South Africa. \textit{SACEMA Research Days (MSc Student Presentation)}, Stellenbosch, South Africa.}

}%
\cvversion{detailed,academic,epidemiologist,datascientist}{%
\needspace{6\baselineskip}
\subsection{Invited Talks}
\cvitem{Nov, 2024}{Epiverse-TRACE: Advancing and Enabling Collaborative Open-Source R Tools for Outbreak Analytics, JUNIPER Weekly Seminar (virtual)}

}%
\cvversion{detailed,academic,academic_twopager,epidemiologist,epidemiologist_twopager}{%
\cvitem{Nov, 2024}{Summary of analytical and statistical tools available for modelling pandemic H5N1 containment and mitigation, WHO R\&D Blueprint Virtual Consultation}

}%
\cvversion{detailed,academic,epidemiologist,datascientist,epidemiologist_twopager}{%
\cvitem{Sep, 2024}{Mpox 2024 outbreak: Common analytics tasks and available R tools/packages. Presented to the UK Pandemic Sciences Network and the World Health Organization Mpox Analytics Collaboratory Seminars}

}%
\cvversion{detailed,academic,epidemiologist}{%
\cvitem{Jun, 2022}{\textit{Modelling outbreak response strategies for decision-making} (\textit{Presented virtually}), Center for Epidemiological Modelling and Analysis, University of Nairobi, Kenya.}

\cvitem{Jul, 2020}{\textit{Modelling the benefits of using the measles vaccine outside cold chain for outbreak response}, Médecins Sans Frontières (MSF) Vaccine Working Group meeting (Presented virtually)}

}%
\cvversion{detailed,academic}{%
\cvitem{Apr, 2019}{\textit{Measles outbreak response vaccination: the role of the vaccine supply chain in low-income settings}, Department of Logistics, Stellenbosch University}

\needspace{6\baselineskip}
\subsection{Poster Presentations}
\cvitem{Dec, 2019}{\textbf{JM Azam}, MJ Ferrari, JRC Pulliam. A framework for evaluating vaccine delivery strategies for responding to measles outbreaks in hard-to-reach populations. \textit{Epidemics7, Charleston, United States (Presented on my behalf).}}

\cvitem{Sep, 2018}{\textbf{JM Azam}, JRC Pulliam. Measles outbreak response vaccination: how can we optimise the vaccine supply chain for efficient responses? \textit{SACEMA Research Days (PhD Student Presentation)}, Stellenbosch, South Africa.}

\cvitem{Jun, 2016}{\textbf{JM Azam}, MI Andersson, E du Toit, M Nieuwoudt. Modelling Hepatitis B Transmission in sub-Saharan Africa - A Focus on the Impact of Vaccination Within Existing HIV MTCT Infrastructure. \textit{Meaningful Modelling of Epidemiological Data, MMED (Student Presentation)}, African Institute for Mathematical Sciences (AIMS), Muizenberg, South Africa.}

\needspace{6\baselineskip}
\section{Panel Discussions}
\cvitem{Nov, 2024}{\textbf{WHO R\&D Blueprint Virtual Consultation.} \textit{Preparing for containment and mitigation of pandemic H5N1 influenza.} We discussed gaps in our knowledge and present approaches to tackling a potential H5N1 widespread outbreak.}

\cvitem{Jul, 2023}{\textbf{Health Data Research UK Black Internship Programme.} \textit{Career Panel.} We shared our career journeys and experiences, and answered questions from the mentees.}

\cvitem{Dec, 2020}{\textbf{The New African Scientist.} \textit{STEM - For Us, By Us conference, Visibility STEM Africa.} We discussed topics including the future of the academic/research environment in Science Technology Engineering Mathematics (STEM) in Africa, mental health among young African researchers, how we find the motivation to do work outside of research, and our vision for STEM in Africa in the next few years.}

\cvitem{Dec, 2018}{\textbf{Science Communication for Policy Engagement.} \textit{Science Forum South Africa (SFSA), CSIR International Convention Centre, Pretoria, South Africa.} We discussed the relationship between science communication and science engagement, and the challenges that early career researchers face in communicating the policy implications of their research. We also discussed (South) Africa's readiness for the Fourth Industrial Revolution and the emerging challenges.}

}%
\needspace{6\baselineskip}
\section{Teaching Experience}
\ifcvversion{academic_twopager,epidemiologist_twopager,datascientist_twopager}{%
\cvversion{academic_twopager,datascientist_twopager}{\cventry{Nov, 2026 -- Present}{Facilitator and tutor}{Data Challenge module, MSc Health Data Science}{London School of Hygiene \& Tropical Medicine}{London, United Kingdom}{}}%
\cvversion{academic_twopager,epidemiologist_twopager}{\cventry{2024 -- 2026}{Lead facilitator}{Modelling for Pandemic Preparedness and Response}{German-West African Centre for Global Health and Pandemic Prevention (G-WAC)}{Kwame Nkrumah University of Science and Technology, Ghana}{Three iterations of the two-week short course, with over 65 participants in total from more than 15 African countries.}}%
\cventry{Jul, 2025 \& Jul, 2026}{Tutor}{Outbreak Analytics and Applied Modelling in R}{London School of Hygiene \& Tropical Medicine}{London, United Kingdom}{}
\cvversion{epidemiologist_twopager}{\cventry{Jun, 2025}{Co-lead facilitator}{Outbreak Data Science for Lab Scientists: A Hands-On Introduction with the Epiverse R Package Ecosystem}{KEMRI Wellcome Trust Research Programme}{Kilifi, Kenya}{}}%
}{%
\cventry{Nov, 2026 -- Present}{Facilitator and tutor}{Data Challenge module, MSc Health Data Science}{London School of Hygiene \& Tropical Medicine}{London, United Kingdom}{Facilitate the Git and GitHub, and R for data exploration and visualisation sessions; tutor a group of 4--5 students on data challenges set by external organisations, including UKHSA, UNICEF, private start-ups, and public health data science companies.}
\cventry{2024 -- 2026}{Lead facilitator}{Modelling for Pandemic Preparedness and Response}{German-West African Centre for Global Health and Pandemic Prevention (G-WAC)}{Kwame Nkrumah University of Science and Technology, Ghana}{Three iterations of the two-week short course, with over 65 participants in total from more than 15 African countries, and an overall course rating above 4.2 out of 5 from participants in every iteration. Coordinate the scientific content and the team of facilitators and tutors; teach the sessions on introduction to compartmental modelling and on model fitting and calibration; and supervise participants' one-month capstone projects.}
\cventry{Jul, 2025 \& Jul, 2026}{Tutor}{Outbreak Analytics and Applied Modelling in R}{London School of Hygiene \& Tropical Medicine}{London, United Kingdom}{Tutored the sessions on accessing epidemiological delays to estimate transmission and severity using the R packages \texttt{epiparameter}, \texttt{EpiNow2}, and \texttt{cfr}. Supported participants in applying the course materials to their own analytical problems and answered their questions.}
\cventry{Jun, 2025}{Co-lead facilitator}{Outbreak Data Science for Lab Scientists: A Hands-On Introduction with the Epiverse R Package Ecosystem}{KEMRI Wellcome Trust Research Programme}{Kilifi, Kenya}{}
\cventry{Apr, 2025}{Co-facilitator}{Hackathon on software tools for rapid outbreak analytics during humanitarian crises}{Humanitarian Modelling Dialogue Workshop, WHO Pandemic Hub}{Berlin, Germany}{}

\cventry{Aug, 2022}{Co-lead facilitator}{Introduction to R: Management, Exploration, and Communication of Data}{SACEMA}{South Africa}{}

}%
\cvversion{detailed,academic}{%
\cventry{Oct, 2017; May, 2018 \& 2019}{Tutor (assistant)}{Introduction to R: Management, Exploration, and Communication of Data}{SACEMA}{South Africa}{}

\cventry{Aug, 2016 -- Jun, 2017}{Teaching assistant}{Introductory Calculus and Algebra}{Department of Mathematical Sciences, Stellenbosch University}{South Africa}{}

\cventry{Sep, 2013 -- Jul, 2014}{Teaching assistant}{Introductory Probability and Statistics, and R for Statistical Analyses}{Department of Mathematics, Kwame Nkrumah University of Science and Technology}{Ghana}{}

}%
\cvversion{academic_twopager,epidemiologist_twopager,datascientist_twopager}{%
\needspace{6\baselineskip}
\section{Student Supervision and Mentoring}
\cvversion{academic_twopager,epidemiologist_twopager}{\cventry{Apr, 2026 -- Apr, 2027}{Mentor}{Measles Analytics Hub Mentoring Scheme}{Imperial College London}{United Kingdom}{}}%
\cventry{Jul, 2025 -- Aug, 2025}{Supervisor}{Health Data Research UK Black Internship Programme}{London School of Hygiene \& Tropical Medicine}{United Kingdom}{Supervised an 8-week project assessing the performance of COVID-19 forecasts in the UK using \texttt{R}, \texttt{EpiNow2}, and \texttt{scoringutils}.}
}%
\cvversion{detailed,academic,epidemiologist,datascientist}{%
\needspace{10\baselineskip}
\section{Student Supervision and Mentoring}
\cventry{Apr, 2026 -- Apr, 2027}{Mentor}{Measles Analytics Hub Mentoring Scheme}{Imperial College London}{United Kingdom}{Mentoring and supervising a project on estimating subnational measles immunity in Uganda. The scheme builds analytics capacity in priority geographies by pairing early- and mid-career analysts with experienced mentors, who provide tailored guidance on modelling methods, epidemiological interpretation, data curation, reproducible workflows, and translating analyses into actionable insights for national immunisation programmes.}
\cventry{Jul, 2025 -- Aug, 2025}{Supervisor}{Health Data Research UK Black Internship Programme}{London School of Hygiene \& Tropical Medicine}{United Kingdom}{I supervised an intern over an 8-week period to develop applied skills in infectious disease forecasting and evaluation through a project assessing the performance of COVID-19 forecasts in the UK using \texttt{R}, \texttt{EpiNow2}, and \texttt{scoringutils}.}
\cventry{Jul, 2024 -- Aug, 2024}{Mentor}{Health Data Research UK Black Internship Programme}{London School of Hygiene \& Tropical Medicine}{London, United Kingdom}{For 8 weeks, I mentored an intern by providing advice on career trajectories, how to effectively conduct his research, and personal development in preparation for PhD studies.}
\cventry{Apr, 2023 -- Jun, 2023}{Mentor}{Research Software Camp}{The Software Sustainability Institute}{United Kingdom}{In 3 months, I mentored a PhD student to develop a working knowledge of using R for running generalised linear mixed models. We worked our way up from no knowledge of R to running and interpreting regression models.}

\needspace{10\baselineskip}
\section{Professional Development}

}%
\cvversion{detailed,academic,epidemiologist}{%
\needspace{6\baselineskip}
\subsection{Leadership and career development}
\cventry{Dec, 2025 -- Feb, 2026}{Diversifying Leadership Programme}{Advance HE}{London}{United Kingdom}{Three-month leadership development programme. Developed skills in navigating institutional power dynamics through cultural capital and strategic networking, in using Action Learning Sets to solve complex leadership challenges, and in articulating a personal leadership story.}
\cvversion{detailed,academic}{\cventry{Apr, 2024 -- Apr, 2025}{B-Mentor Mentoring Programme}{University College London}{London}{United Kingdom}{Year-long mentoring from a Senior Research Fellow at UCL on moving from research software engineering into a research fellow career, including identifying grant and fellowship opportunities, building academic networks, and building a distinctive academic profile. My mentor reviewed the fellowship proposal I developed on LSHTM's Grant Application Programme.}}

\needspace{6\baselineskip}
\subsection{Grant writing}
\cventry{Jan, 2025 -- Jun, 2025}{Grant Application Programme}{London School of Hygiene \& Tropical Medicine}{London}{United Kingdom}{Selected for LSHTM's six-month programme, run by the Strategic Research Office, to develop future research leaders by building skills in research proposal development. Covered choosing funding schemes and timing applications, writing for funding committees and expert referees, and critically reviewing applications, through workshops, panel discussions, mock grant and fellowship review boards, and one-to-one mentorship from an LSHTM Professor. Developed a proposal, \textit{Improving Infectious Disease Forecasting Methodology and Capacity Development for Public Health Impact in Africa}, for submission to a future funding call.}

}%
\cvversion{detailed,academic}{%
\needspace{6\baselineskip}
\subsection{Teaching and curriculum development}
\cventry{Dec, 2023}{Curriculum development workshop for the short course, ``Modelling for Pandemic Preparedness and Response''}{German-West African Centre for Global Health and Pandemic Prevention/Technische Universität Berlin}{Berlin}{Germany}{This workshop involved a meeting of experts to design the named two week short course to be offered in Ghana.}

}%
\cvversion{detailed,datascientist}{%
\needspace{6\baselineskip}
\subsection{Technical and methods training}
\cventry{Jan, 2020}{Software Engineering for Applied Mathematical Sciences}{SACEMA}{Stellenbosch}{}{}
}%
\cvversion{detailed}{%
\cventry{Nov, 2019}{The Annual African Vaccinology Course}{Vaccines for Africa Initiative (VACFA)}{University of Cape Town, South Africa}{}{}
\cventry{Jul, 2018}{The Systematic Review}{African Doctoral Academy}{Stellenbosch University}{South Africa}{}
\cventry{May, 2018}{Individual-based Modelling in Epidemiology - A Practical Introduction}{SACEMA}{Stellenbosch University}{South Africa}{}
}%
\cvversion{detailed,datascientist}{%
\cventry{Dec, 2017}{Spatial Disease Mapping}{University of Cape Town}{Cape Town}{}{}
}%
\cvversion{detailed}{%
\cventry{Aug, 2017}{Introduction to Scientific Writing Skills}{The Writing Center}{Stellenbosch University}{South Africa}{}
}%
\cvversion{detailed,datascientist}{%
\cventry{May, 2017}{Data Science}{AIMS-South Africa}{Muizenberg, Cape Town, South Africa}{}{}
\cventry{Oct, 2016}{Bayesian Analysis of Longitudinal Data}{SACEMA}{Stellenbosch University}{South Africa}{}
}%
\cvversion{detailed}{%
\cventry{Sep, 2016}{Mathematical Modeling of Infectious Diseases (MMID)}{University of Cape Town}{Cape Town, South Africa}{}{}
\cventry{Jul, 2016}{Clinic on the Meaningful Modeling of Epidemiological Data (MMED)}{University of Florida, SACEMA-Stellenbosch University, AIMS-South Africa}{Muizenberg, Cape Town, South Africa}{}{}
\cventry{May, 2016}{Applying Quantitative Bias Analysis to Epidemiologic Data}{SACEMA}{Stellenbosch University}{South Africa}{}
\cventry{Mar, 2016 -- Aug, 2016}{Introduction to Epidemiology and Research Methodology}{SACEMA}{Stellenbosch University}{South Africa}{}
}%
\cvversion{detailed,datascientist}{%
\cventry{Mar, 2016}{Bayesian Biostatistics}{SACEMA}{Stellenbosch University}{South Africa}{}
}%
\cvversion{detailed}{%
\cventry{Jul, 2015}{Computing in Applied Mathematical Sciences (CAMS)}
{AIMS-Ghana}{}{}{}
}%

\cvversion{detailed,epidemiologist,datascientist}{%
\needspace{6\baselineskip}
\subsection{Meetings and workshops}
}%
\cvversion{detailed,epidemiologist}{%
\cventry{Feb, 2025}{UK Epidemic Analytics and Modelling Meeting}{London School of Hygiene \& Tropical Medicine (HPRU in Health Analytics and Modelling)}{London}{United Kingdom}{This was a 2-day meeting of modellers from across the UK to network and brainstorm the way forward with pro-active preparedness towards a more efficient epidemic analytics in the UK.}
}%
\cvversion{detailed,epidemiologist,datascientist}{%
\cventry{Nov, 2024}{Understanding, evaluating, and improving forecasts of infectious disease burden}{Infectious Disease Modelling Conference}{Bangkok}{Thailand}{This was a half-day pre-conference workshop.}
}%
\cvversion{detailed,datascientist}{%
\cventry{Apr, 2024}{R Developer Day}{Imperial College London}{London}{United Kingdom}{An open, collaborative event for people interested in contributing to base R. I worked with two others to submit a patch for an improvement to the base package's substr() and substring() functions to help clarify their similarities and differences and demonstrate some previously obscure features.}
\cventry{Jul, 2023}{RSLondonSouthEast 2023}{Imperial College London}{London}{United Kingdom}{This workshop aimed to provide a regional focus on research software activities, offer an opportunity to present work to a regional audience, and help participants develop links with the wider Research Software Engineering community.}
\cventry{Mar, 2023}{Forecasting natural and social systems}{The Royal Society}{London}{United Kingdom}{This meeting aimed to convene leading experts in forecasting across disciplines to discuss the predictability of different systems, to exchange expertise, and to identify common challenges and possible solutions.}
}%
\cvversion{detailed}{%
\cventry{Dec, 2022}{100 days and 100 lines of code}{Epiverse-TRACE, London School of Hygiene \& Tropical Medicine}{London}{United Kingdom}{}
}%
\cvversion{detailed,epidemiologist}{%
\cventry{Oct, 2022}{Measles modelling workshop}{Vaccine Impact Modelling Consortium}{Virtual}{}{}
\cventry{Nov, 2017}{The Use of Measles Vaccines Under Extended Controlled Temperature Conditions}{Médecins Sans Frontières}{Paris, France}{}{}
}%
\cvversion{detailed}{%
\cventry{May, 2016}{Global Change Impact on Infectious Diseases and Alien Species Expansion}
{AIMS-South Africa in collaboration with 16 other societies/organisations}{}{}{}

\needspace{6\baselineskip}
\subsection{Science communication}
\cventry{Nov, 2020}{FameLab - Science Communication Competition}{The South African Centre for Epidemiological Modelling and Analysis (SACEMA)}{Stellenbosch, South Africa}{}{}
\cventry{Nov, 2020}{Learn to write for the media}{The Conversation Africa}{}{}{An online course that teaches the basics of how to write media/popular articles based on scientific findings.}
\cventry{Dec, 2019}{Science Communication Master Class}{Stellenbosch University}{Stellenbosch}{}{}
\cventry{Feb, 2019}{FameLab - Science Communication Competition}{The South African Centre for Epidemiological Modelling and Analysis (SACEMA)}{Stellenbosch, South Africa}{}{}
\cventry{May, 2018}{The Conversation Writing Workshop}{Stellenbosch University}{South Africa}{}{}
\cventry{Dec, 2017}{FameLab - Science Communication Competition}{CSIR International Convention Centre}{Pretoria, South Africa}{}{}
\cventry{Aug, 2016}{Science Communication and Public Engagement}{AIMS-South Africa}{Muizenberg, Cape Town, South Africa}{}{}

}%
% \needspace{6\baselineskip}
% \subsection{Equality and compliance training}
% \cventry{Sep, 2023}{Equality \& Diversity Essentials}{London School of Hygiene \& Tropical Medicine}{London}{}{}
% \cventry{Nov, 2022}{GDPR Compliance: Essential Training}{LinkedIn Learning}{}{}{}

\cvversion{detailed,academic,epidemiologist}{%
\needspace{6\baselineskip}
\section{Technical Competencies}
\cvtechnicalcompetencies
}%
\cvversion{academic_twopager,epidemiologist_twopager}{%
\needspace{6\baselineskip}
\section{Technical Competencies}
\cvitem{Programming tools}{R, Stan, Bash, GNU Make, GitHub \& Git}
\cvitem{Machine learning}{Tree-based methods and gradient boosting, using \texttt{tidymodels}}
}%

\cvversion{detailed,academic,epidemiologist,datascientist}{%
\needspace{6\baselineskip}
\section{Professional Memberships}
\cventry{Sep, 2024 -- Present}{Affiliate Member}{Joint UNIversities Pandemic and Epidemiological Research (JUNIPER)}{United Kingdom}{}{}
\cventry{Sep, 2023 -- Present}{Member}{Models of Infectious Disease Agent Study (MIDAS)}{University of Pittsburgh}{United States}{}
\cventry{Dec, 2022 -- Present}{Member}{Society of Research Software Engineering}{United Kingdom}{}{}

}%
\cvversion{detailed,academic,epidemiologist,datascientist}{%
\needspace{6\baselineskip}
\section{Referees}
\cvitem{}{Available upon request.}
}%

\end{document}
